%% ============================================================
%% COMM514 MSc Report -- Section 2: Literature Review
%% University of Exeter
%% ============================================================
%% Embed in your main document with %% ============================================================
%% COMM514 MSc Report -- Section 2: Literature Review
%% University of Exeter
%% ============================================================
%% Embed in your main document with %% ============================================================
%% COMM514 MSc Report -- Section 2: Literature Review
%% University of Exeter
%% ============================================================
%% Embed in your main document with %% ============================================================
%% COMM514 MSc Report -- Section 2: Literature Review
%% University of Exeter
%% ============================================================
%% Embed in your main document with \input{literature_review}
%% Requires the same preamble as introduction.tex
%% ============================================================

\section{Literature Review}

\subsection{AI-Assisted Programming}

The emergence of large language models trained on code has transformed the ambition of automated programming from the generation of short snippets to the synthesis of complete functional programs. \citet{chen2021} introduced the HumanEval benchmark and the \textit{pass@k} evaluation metric, establishing functional correctness as the field's primary measure of progress. By framing code generation as a problem of producing programs that pass unit tests, HumanEval implicitly defined capability as the ability to satisfy a functional specification; security properties received no representation in the benchmark. The influence of this design choice on subsequent research has been substantial. \citet{li2022alphacode} evaluated AlphaCode on competitive programming contests, achieving [Li et al. 2022 -- insert percentile ranking], a result celebrated as evidence of near-human-level programming ability. \citet{roziere2023} presented Code Llama, demonstrating [Roziere et al. 2023 -- insert benchmark result on HumanEval or equivalent], again evaluated primarily against correctness benchmarks. Across this lineage of models, the optimisation target has been consistent: produce code that executes and passes tests.

The consequences of this framing have been documented empirically. \citet{liu2023} demonstrated that code generated by large language models and evaluated as functionally correct by standard benchmarks nonetheless contains security vulnerabilities at a measurable rate, revealing a systematic gap between the property being measured and the property that matters in deployment. Correctness and security are not the same objective; a function can pass every test in a suite while simultaneously being exploitable if the test suite does not probe security-relevant inputs. The productivity gains that have driven commercial adoption are real: \citet{peng2023} found that developers using AI coding assistance completed tasks [Peng et al. 2023 -- insert productivity improvement, e.g. percentage faster] faster than those working without it. The practical implication is that AI-assisted code is being shipped at an accelerating rate, under the assumption that faster code generation is an unambiguous benefit, without adequate attention to what properties the generated code does and does not possess. The field has, in effect, solved the wrong problem at scale.

\subsection{Security of AI-Generated Code}

The question of whether AI-generated code is secure has attracted a growing body of empirical work, and the findings converge on a troubling pattern that grows more complex as the body of evidence accumulates. \citet{pearce2022} conducted the first large-scale systematic evaluation, generating 1,689 programs across a range of scenarios drawn from the Common Weakness Enumeration taxonomy and finding that approximately 40 per cent contained exploitable security vulnerabilities. Critically, this vulnerability rate was not limited to a subset of task types; it was distributed across diverse CWE categories, establishing that insecure code generation is a property of the model's general behaviour rather than a domain-specific failure mode.

\citet{khoury2023} extended this finding in a particularly revealing direction by demonstrating that ChatGPT produces insecure code even in circumstances where the model itself, when prompted directly, can identify the security principle being violated. This dissociation between declarative security knowledge and generative security practice suggests that the problem is not one of ignorance but of application: the model possesses relevant knowledge but does not activate it during code generation without explicit instruction to do so.

The most consequential finding in this literature, however, comes from \citet{perry2023}, who examined not just the properties of AI-generated code but the effect of AI assistance on developer behaviour. Their controlled study found that participants using AI assistance wrote less secure code than unassisted participants on four out of five tasks, including a particularly striking result on an ECDSA signing task where the proportion of participants producing a secure implementation dropped from 21 per cent in the unassisted condition to 3 per cent in the AI-assisted condition, a difference that was statistically significant ($p = 0.039$). More troubling still, participants in the AI-assisted condition rated their own code as significantly more secure than those in the control condition, despite their outputs being objectively less secure. The confidence paradox documented by Perry et al. constitutes a distinct and underappreciated dimension of the risk: AI assistance does not merely fail to improve security, it actively degrades the developer's epistemic posture toward security.

The breadth of the vulnerability problem across CWE categories has been characterised more formally by \citet{siddiq2022}, whose SecurityEval benchmark provides 130 security-specific samples spanning 75 CWE types, offering a more structured evaluation instrument than task-based studies permit. Taken together, the findings of Pearce et al., Khoury et al., Perry et al., and Siddiq and Santos trace three distinct but related failure modes: AI generates vulnerable code, AI-assisted developers are less likely to notice or correct that vulnerability, and the developer's confidence in the insecure output is, if anything, elevated.

\citet{shukla2025} introduce a fourth failure mode that carries particular implications for iterative AI-assisted workflows. Their study of LLM code refinement found that repeated cycles of AI feedback and revision produced a 37.6 per cent increase in critical vulnerabilities after five iterations. Rather than converging on more secure outputs through successive refinement, the iterative process amplified security weaknesses. This finding challenges the intuitive assumption that more AI involvement in the development process produces better outcomes and suggests that the relationship between the quantity of AI-generated content and code security is not merely neutral but potentially negative.

\subsection{Multi-Agent Systems and Structured Reasoning}

The limitations of single-agent LLM systems on complex tasks have motivated a parallel body of research into multi-agent architectures, in which the division of cognitive labour across specialised agents is proposed as a mechanism for improving output quality. \citet{li2023camel} introduced the CAMEL framework, demonstrating that role-based decomposition, in which agents adopt defined functional identities within a collaborative workflow, improved performance on tasks requiring extended reasoning relative to single-agent baselines. \citet{wu2023autogen} developed AutoGen, a framework for multi-agent conversational systems that enables flexible orchestration of agent roles, showing that conversation-based multi-agent coordination consistently outperformed single-agent approaches on complex compositional tasks.

The theoretical basis for these improvements was partly anticipated by \citet{yao2023}, whose ReAct framework demonstrated that interleaving explicit reasoning traces with action selection produced more reliable task completion than either pure reasoning or pure action. The relevance of ReAct to agentic security pipelines is structural: if explicit reasoning prior to action improves outcome quality, then a pipeline architecture that mandates explicit threat analysis prior to code generation replicates this mechanism at the level of agent orchestration rather than individual token prediction.

The most directly relevant empirical evidence for the deployment of multi-agent human-in-the-loop systems comes from \citet{takerngsaksiri2025}, who documented the HULA system developed and deployed at Atlassian. Across a production workload, HULA achieved plan success in 79 per cent of cases, code success in 87 per cent of cases, and a 59 per cent pull request merge rate, with 568 tasks terminated early by human reviewers exercising their oversight role. These figures represent the most rigorous published evidence for the real-world viability of human-orchestrated agentic pipelines in a software development context. However, the HULA evaluation focused exclusively on task completion and development efficiency. Security outcomes were not measured, vulnerability density was not reported, and the relationship between human oversight and the security properties of the produced code was not investigated. HULA thus demonstrates that the architectural category is viable and scalable; it does not establish whether that architecture produces more secure code than a simpler alternative.

\subsection{Human-AI Collaboration}

The question of how to design effective human oversight of AI systems predates the current generation of code generation tools and has produced a body of design guidance that the AI-assisted programming literature has been slow to engage with. \citet{amershi2019} synthesised 18 guidelines for human-AI interaction from human-computer interaction research, establishing transparency, controllability, and calibrated confidence as foundational requirements for systems that integrate human and AI decision-making. Calibration is particularly relevant in the context of AI-assisted code generation: Perry et al.'s confidence paradox indicates that current tools violate the calibration principle systematically, producing higher confidence in developers when confidence is least warranted.

\citet{shneiderman2020} challenged the prevailing assumption that high automation and high human control represent opposite ends of a spectrum, arguing instead that effective human-centred AI systems achieve both simultaneously through appropriate interface and workflow design. The implication is that reducing AI automation to increase human control is not the correct response to the risks documented by Perry et al. and Shukla et al.; the correct response is to redesign the workflow so that human oversight is structurally effective rather than structurally bypassed.

The evidence on whether human oversight is structurally effective in practice is not encouraging. \citet{perry2023} observed that existing AI coding tools position the developer primarily as a passive acceptor of suggestions, with the cognitive effort required to evaluate a suggestion for security properties working in tension with the productivity incentive that motivates AI adoption in the first place. \citet{kennedy2025} examined human oversight in AI governance contexts more broadly, finding that humans embedded in oversight loops provided correct oversight only approximately half the time, with failure driven predominantly by institutional and social pressure to approve rather than to scrutinise. This finding is significant because it demonstrates that the mere presence of a human decision point in an AI workflow does not guarantee meaningful oversight; the design of that decision point, including the information provided, the framing of the choice, and the presence or absence of structured evaluation criteria, determines whether oversight is genuine or nominal.

\subsection{Research Gap}

The four bodies of literature reviewed above identify a convergent problem that none of them individually resolves. The AI-assisted programming literature has established that current models optimise for correctness and achieve it at scale, but correctness and security are independent properties. The security literature has established that AI-generated code is insecure at a substantial rate, that developer vigilance is reduced by AI assistance, and that iterative refinement compounds rather than corrects these deficiencies. The multi-agent literature has established that structured role-based architectures improve reasoning quality on complex tasks and has demonstrated their viability in production software development, but has not examined security outcomes. The human-AI collaboration literature has established the design principles for effective oversight but has also documented the tendency of oversight mechanisms to become nominal in practice without structural support for genuine scrutiny.

What remains absent from the empirical record is a controlled experiment that brings these threads together: a study that directly compares a standard single-agent code generation baseline against a structured multi-agent pipeline incorporating pre-generation threat modelling, independent review, and formal verification, with security outcome as the primary dependent variable and developer reasoning quality as an explicit secondary outcome. No existing study has reported vulnerability density as a comparative metric across these two architectural conditions under equivalent task and participant constraints. The present study is designed to fill that gap. By holding the underlying model, temperature, and task set constant across conditions, and varying only the architectural structure through which code generation proceeds, the study isolates the effect of multi-agent orchestration on security outcomes in a way that existing work does not permit. The explicit measurement of participant decisions at human-in-the-loop checkpoints further distinguishes this study from both the single-agent vulnerability literature and the multi-agent task-completion literature by treating developer security reasoning as an outcome in its own right rather than an incidental observation.

%% ============================================================
%% BIBLIOGRAPHY ENTRIES TO ADD TO YOUR .bib FILE
%% Add these to the existing entries from introduction.tex
%% ============================================================

%% \bibitem[Amershi et al.(2019)]{amershi2019}
%% Amershi, S., Weld, D., Vorvoreanu, M., Fourney, A., Nushi, B.,
%% Collisson, P., Suh, J., Iqbal, S., Bennett, P., Inkpen, K., Teevan, J.,
%% Kikin-Gil, R. and Horvitz, E., 2019.
%% Software engineering for machine learning: A case study.
%% OR: Guidelines for Human-AI Interaction [verify full title from source].
%% \newblock In: \emph{Proceedings of the 2019 CHI Conference on Human Factors
%% in Computing Systems (CHI '19)}, Glasgow, UK. ACM, pp.~[insert pages].

%% \bibitem[Chen et al.(2021)]{chen2021}
%% Chen, M., Tworek, J., Jun, H., Yuan, Q., Pinto, H.P.O., Kaplan, J., ...
%% and Zaremba, W., 2021.
%% \newblock Evaluating large language models trained on code.
%% \newblock \emph{arXiv preprint arXiv:2107.03374}.

%% \bibitem[Khoury et al.(2023)]{khoury2023}
%% Khoury, R., Avila, A.R., Brunelle, J. and Camara, B.M., 2023.
%% \newblock How secure is code generated by ChatGPT?
%% \newblock In: \emph{[Conference -- insert from paper]}, [Location].
%% IEEE/ACM, pp.~[insert pages].

%% \bibitem[Li et al.(2022)]{li2022alphacode}
%% Li, Y., Choi, D., Chung, J., Kushman, N., Schrittwieser, J., Leblond, R.,
%% ... and Vinyals, O., 2022.
%% \newblock Competition-level code generation with AlphaCode.
%% \newblock \emph{Science}, 378(6624), pp.~1092--1097.

%% \bibitem[Li et al.(2023)]{li2023camel}
%% Li, G., Hammoud, H.A.A.K., Itani, H., Khizbullin, D. and
%% Ghanem, B., 2023.
%% \newblock CAMEL: Communicative agents for ``mind'' exploration of large
%% language model society.
%% \newblock In: \emph{Advances in Neural Information Processing Systems (NeurIPS 2023)},
%% [Location]. NeurIPS, pp.~[insert pages].

%% \bibitem[Liu et al.(2023)]{liu2023}
%% Liu, J. et al., 2023.
%% \newblock [Full title -- verify from paper; relates to security errors in
%% functionally correct AI-generated code].
%% \newblock \emph{[Venue -- insert from paper]}, pp.~[pages].

%% \bibitem[Peng et al.(2023)]{peng2023}
%% Peng, S., Kalliamvakou, E., Croft, P. and Lombard, C., 2023.
%% \newblock The impact of AI on developer productivity: Evidence from
%% GitHub Copilot.
%% \newblock \emph{arXiv preprint arXiv:2302.06590}.

%% \bibitem[Roziere et al.(2023)]{roziere2023}
%% Roziere, B., Gehring, J., Gloeckle, F., Sootla, S., Gat, I., Alon, U.,
%% ... and Synnaeve, G., 2023.
%% \newblock Code Llama: Open foundation models for code.
%% \newblock \emph{arXiv preprint arXiv:2308.12950}.

%% \bibitem[Shneiderman(2020)]{shneiderman2020}
%% Shneiderman, B., 2020.
%% \newblock Human-centered artificial intelligence: Three fresh ideas.
%% \newblock \emph{AIS Transactions on Human-Computer Interaction}, 12(3),
%% pp.~109--124.

%% \bibitem[Siddiq and Santos(2022)]{siddiq2022}
%% Siddiq, M.L. and Santos, J.C.S., 2022.
%% \newblock SecurityEval dataset: Mining vulnerability examples to evaluate
%% machine learning-based code generation techniques.
%% \newblock In: \emph{Proceedings of the 4th International Workshop on
%% Mining Software Repositories for Privacy and Security (MSR4P\&S '22)}.
%% ACM, pp.~[insert pages].

%% \bibitem[Wu et al.(2023)]{wu2023autogen}
%% Wu, Q., Bansal, G., Zhang, J., Wu, Y., Li, B., Zhu, E., ... and
%% Wang, C., 2023.
%% \newblock AutoGen: Enabling next-generation LLM applications via
%% multi-agent conversation.
%% \newblock \emph{arXiv preprint arXiv:2308.08155}.

%% \bibitem[Yao et al.(2023)]{yao2023}
%% Yao, S., Zhao, J., Yu, D., Du, N., Shafran, I., Narasimhan, K.
%% and Cao, Y., 2023.
%% \newblock ReAct: Synergizing reasoning and acting in language models.
%% \newblock In: \emph{Proceedings of the 11th International Conference on
%% Learning Representations (ICLR 2023)}.

%% Requires the same preamble as introduction.tex
%% ============================================================

\section{Literature Review}

\subsection{AI-Assisted Programming}

The emergence of large language models trained on code has transformed the ambition of automated programming from the generation of short snippets to the synthesis of complete functional programs. \citet{chen2021} introduced the HumanEval benchmark and the \textit{pass@k} evaluation metric, establishing functional correctness as the field's primary measure of progress. By framing code generation as a problem of producing programs that pass unit tests, HumanEval implicitly defined capability as the ability to satisfy a functional specification; security properties received no representation in the benchmark. The influence of this design choice on subsequent research has been substantial. \citet{li2022alphacode} evaluated AlphaCode on competitive programming contests, achieving [Li et al. 2022 -- insert percentile ranking], a result celebrated as evidence of near-human-level programming ability. \citet{roziere2023} presented Code Llama, demonstrating [Roziere et al. 2023 -- insert benchmark result on HumanEval or equivalent], again evaluated primarily against correctness benchmarks. Across this lineage of models, the optimisation target has been consistent: produce code that executes and passes tests.

The consequences of this framing have been documented empirically. \citet{liu2023} demonstrated that code generated by large language models and evaluated as functionally correct by standard benchmarks nonetheless contains security vulnerabilities at a measurable rate, revealing a systematic gap between the property being measured and the property that matters in deployment. Correctness and security are not the same objective; a function can pass every test in a suite while simultaneously being exploitable if the test suite does not probe security-relevant inputs. The productivity gains that have driven commercial adoption are real: \citet{peng2023} found that developers using AI coding assistance completed tasks [Peng et al. 2023 -- insert productivity improvement, e.g. percentage faster] faster than those working without it. The practical implication is that AI-assisted code is being shipped at an accelerating rate, under the assumption that faster code generation is an unambiguous benefit, without adequate attention to what properties the generated code does and does not possess. The field has, in effect, solved the wrong problem at scale.

\subsection{Security of AI-Generated Code}

The question of whether AI-generated code is secure has attracted a growing body of empirical work, and the findings converge on a troubling pattern that grows more complex as the body of evidence accumulates. \citet{pearce2022} conducted the first large-scale systematic evaluation, generating 1,689 programs across a range of scenarios drawn from the Common Weakness Enumeration taxonomy and finding that approximately 40 per cent contained exploitable security vulnerabilities. Critically, this vulnerability rate was not limited to a subset of task types; it was distributed across diverse CWE categories, establishing that insecure code generation is a property of the model's general behaviour rather than a domain-specific failure mode.

\citet{khoury2023} extended this finding in a particularly revealing direction by demonstrating that ChatGPT produces insecure code even in circumstances where the model itself, when prompted directly, can identify the security principle being violated. This dissociation between declarative security knowledge and generative security practice suggests that the problem is not one of ignorance but of application: the model possesses relevant knowledge but does not activate it during code generation without explicit instruction to do so.

The most consequential finding in this literature, however, comes from \citet{perry2023}, who examined not just the properties of AI-generated code but the effect of AI assistance on developer behaviour. Their controlled study found that participants using AI assistance wrote less secure code than unassisted participants on four out of five tasks, including a particularly striking result on an ECDSA signing task where the proportion of participants producing a secure implementation dropped from 21 per cent in the unassisted condition to 3 per cent in the AI-assisted condition, a difference that was statistically significant ($p = 0.039$). More troubling still, participants in the AI-assisted condition rated their own code as significantly more secure than those in the control condition, despite their outputs being objectively less secure. The confidence paradox documented by Perry et al. constitutes a distinct and underappreciated dimension of the risk: AI assistance does not merely fail to improve security, it actively degrades the developer's epistemic posture toward security.

The breadth of the vulnerability problem across CWE categories has been characterised more formally by \citet{siddiq2022}, whose SecurityEval benchmark provides 130 security-specific samples spanning 75 CWE types, offering a more structured evaluation instrument than task-based studies permit. Taken together, the findings of Pearce et al., Khoury et al., Perry et al., and Siddiq and Santos trace three distinct but related failure modes: AI generates vulnerable code, AI-assisted developers are less likely to notice or correct that vulnerability, and the developer's confidence in the insecure output is, if anything, elevated.

\citet{shukla2025} introduce a fourth failure mode that carries particular implications for iterative AI-assisted workflows. Their study of LLM code refinement found that repeated cycles of AI feedback and revision produced a 37.6 per cent increase in critical vulnerabilities after five iterations. Rather than converging on more secure outputs through successive refinement, the iterative process amplified security weaknesses. This finding challenges the intuitive assumption that more AI involvement in the development process produces better outcomes and suggests that the relationship between the quantity of AI-generated content and code security is not merely neutral but potentially negative.

\subsection{Multi-Agent Systems and Structured Reasoning}

The limitations of single-agent LLM systems on complex tasks have motivated a parallel body of research into multi-agent architectures, in which the division of cognitive labour across specialised agents is proposed as a mechanism for improving output quality. \citet{li2023camel} introduced the CAMEL framework, demonstrating that role-based decomposition, in which agents adopt defined functional identities within a collaborative workflow, improved performance on tasks requiring extended reasoning relative to single-agent baselines. \citet{wu2023autogen} developed AutoGen, a framework for multi-agent conversational systems that enables flexible orchestration of agent roles, showing that conversation-based multi-agent coordination consistently outperformed single-agent approaches on complex compositional tasks.

The theoretical basis for these improvements was partly anticipated by \citet{yao2023}, whose ReAct framework demonstrated that interleaving explicit reasoning traces with action selection produced more reliable task completion than either pure reasoning or pure action. The relevance of ReAct to agentic security pipelines is structural: if explicit reasoning prior to action improves outcome quality, then a pipeline architecture that mandates explicit threat analysis prior to code generation replicates this mechanism at the level of agent orchestration rather than individual token prediction.

The most directly relevant empirical evidence for the deployment of multi-agent human-in-the-loop systems comes from \citet{takerngsaksiri2025}, who documented the HULA system developed and deployed at Atlassian. Across a production workload, HULA achieved plan success in 79 per cent of cases, code success in 87 per cent of cases, and a 59 per cent pull request merge rate, with 568 tasks terminated early by human reviewers exercising their oversight role. These figures represent the most rigorous published evidence for the real-world viability of human-orchestrated agentic pipelines in a software development context. However, the HULA evaluation focused exclusively on task completion and development efficiency. Security outcomes were not measured, vulnerability density was not reported, and the relationship between human oversight and the security properties of the produced code was not investigated. HULA thus demonstrates that the architectural category is viable and scalable; it does not establish whether that architecture produces more secure code than a simpler alternative.

\subsection{Human-AI Collaboration}

The question of how to design effective human oversight of AI systems predates the current generation of code generation tools and has produced a body of design guidance that the AI-assisted programming literature has been slow to engage with. \citet{amershi2019} synthesised 18 guidelines for human-AI interaction from human-computer interaction research, establishing transparency, controllability, and calibrated confidence as foundational requirements for systems that integrate human and AI decision-making. Calibration is particularly relevant in the context of AI-assisted code generation: Perry et al.'s confidence paradox indicates that current tools violate the calibration principle systematically, producing higher confidence in developers when confidence is least warranted.

\citet{shneiderman2020} challenged the prevailing assumption that high automation and high human control represent opposite ends of a spectrum, arguing instead that effective human-centred AI systems achieve both simultaneously through appropriate interface and workflow design. The implication is that reducing AI automation to increase human control is not the correct response to the risks documented by Perry et al. and Shukla et al.; the correct response is to redesign the workflow so that human oversight is structurally effective rather than structurally bypassed.

The evidence on whether human oversight is structurally effective in practice is not encouraging. \citet{perry2023} observed that existing AI coding tools position the developer primarily as a passive acceptor of suggestions, with the cognitive effort required to evaluate a suggestion for security properties working in tension with the productivity incentive that motivates AI adoption in the first place. \citet{kennedy2025} examined human oversight in AI governance contexts more broadly, finding that humans embedded in oversight loops provided correct oversight only approximately half the time, with failure driven predominantly by institutional and social pressure to approve rather than to scrutinise. This finding is significant because it demonstrates that the mere presence of a human decision point in an AI workflow does not guarantee meaningful oversight; the design of that decision point, including the information provided, the framing of the choice, and the presence or absence of structured evaluation criteria, determines whether oversight is genuine or nominal.

\subsection{Research Gap}

The four bodies of literature reviewed above identify a convergent problem that none of them individually resolves. The AI-assisted programming literature has established that current models optimise for correctness and achieve it at scale, but correctness and security are independent properties. The security literature has established that AI-generated code is insecure at a substantial rate, that developer vigilance is reduced by AI assistance, and that iterative refinement compounds rather than corrects these deficiencies. The multi-agent literature has established that structured role-based architectures improve reasoning quality on complex tasks and has demonstrated their viability in production software development, but has not examined security outcomes. The human-AI collaboration literature has established the design principles for effective oversight but has also documented the tendency of oversight mechanisms to become nominal in practice without structural support for genuine scrutiny.

What remains absent from the empirical record is a controlled experiment that brings these threads together: a study that directly compares a standard single-agent code generation baseline against a structured multi-agent pipeline incorporating pre-generation threat modelling, independent review, and formal verification, with security outcome as the primary dependent variable and developer reasoning quality as an explicit secondary outcome. No existing study has reported vulnerability density as a comparative metric across these two architectural conditions under equivalent task and participant constraints. The present study is designed to fill that gap. By holding the underlying model, temperature, and task set constant across conditions, and varying only the architectural structure through which code generation proceeds, the study isolates the effect of multi-agent orchestration on security outcomes in a way that existing work does not permit. The explicit measurement of participant decisions at human-in-the-loop checkpoints further distinguishes this study from both the single-agent vulnerability literature and the multi-agent task-completion literature by treating developer security reasoning as an outcome in its own right rather than an incidental observation.

%% ============================================================
%% BIBLIOGRAPHY ENTRIES TO ADD TO YOUR .bib FILE
%% Add these to the existing entries from introduction.tex
%% ============================================================

%% \bibitem[Amershi et al.(2019)]{amershi2019}
%% Amershi, S., Weld, D., Vorvoreanu, M., Fourney, A., Nushi, B.,
%% Collisson, P., Suh, J., Iqbal, S., Bennett, P., Inkpen, K., Teevan, J.,
%% Kikin-Gil, R. and Horvitz, E., 2019.
%% Software engineering for machine learning: A case study.
%% OR: Guidelines for Human-AI Interaction [verify full title from source].
%% \newblock In: \emph{Proceedings of the 2019 CHI Conference on Human Factors
%% in Computing Systems (CHI '19)}, Glasgow, UK. ACM, pp.~[insert pages].

%% \bibitem[Chen et al.(2021)]{chen2021}
%% Chen, M., Tworek, J., Jun, H., Yuan, Q., Pinto, H.P.O., Kaplan, J., ...
%% and Zaremba, W., 2021.
%% \newblock Evaluating large language models trained on code.
%% \newblock \emph{arXiv preprint arXiv:2107.03374}.

%% \bibitem[Khoury et al.(2023)]{khoury2023}
%% Khoury, R., Avila, A.R., Brunelle, J. and Camara, B.M., 2023.
%% \newblock How secure is code generated by ChatGPT?
%% \newblock In: \emph{[Conference -- insert from paper]}, [Location].
%% IEEE/ACM, pp.~[insert pages].

%% \bibitem[Li et al.(2022)]{li2022alphacode}
%% Li, Y., Choi, D., Chung, J., Kushman, N., Schrittwieser, J., Leblond, R.,
%% ... and Vinyals, O., 2022.
%% \newblock Competition-level code generation with AlphaCode.
%% \newblock \emph{Science}, 378(6624), pp.~1092--1097.

%% \bibitem[Li et al.(2023)]{li2023camel}
%% Li, G., Hammoud, H.A.A.K., Itani, H., Khizbullin, D. and
%% Ghanem, B., 2023.
%% \newblock CAMEL: Communicative agents for ``mind'' exploration of large
%% language model society.
%% \newblock In: \emph{Advances in Neural Information Processing Systems (NeurIPS 2023)},
%% [Location]. NeurIPS, pp.~[insert pages].

%% \bibitem[Liu et al.(2023)]{liu2023}
%% Liu, J. et al., 2023.
%% \newblock [Full title -- verify from paper; relates to security errors in
%% functionally correct AI-generated code].
%% \newblock \emph{[Venue -- insert from paper]}, pp.~[pages].

%% \bibitem[Peng et al.(2023)]{peng2023}
%% Peng, S., Kalliamvakou, E., Croft, P. and Lombard, C., 2023.
%% \newblock The impact of AI on developer productivity: Evidence from
%% GitHub Copilot.
%% \newblock \emph{arXiv preprint arXiv:2302.06590}.

%% \bibitem[Roziere et al.(2023)]{roziere2023}
%% Roziere, B., Gehring, J., Gloeckle, F., Sootla, S., Gat, I., Alon, U.,
%% ... and Synnaeve, G., 2023.
%% \newblock Code Llama: Open foundation models for code.
%% \newblock \emph{arXiv preprint arXiv:2308.12950}.

%% \bibitem[Shneiderman(2020)]{shneiderman2020}
%% Shneiderman, B., 2020.
%% \newblock Human-centered artificial intelligence: Three fresh ideas.
%% \newblock \emph{AIS Transactions on Human-Computer Interaction}, 12(3),
%% pp.~109--124.

%% \bibitem[Siddiq and Santos(2022)]{siddiq2022}
%% Siddiq, M.L. and Santos, J.C.S., 2022.
%% \newblock SecurityEval dataset: Mining vulnerability examples to evaluate
%% machine learning-based code generation techniques.
%% \newblock In: \emph{Proceedings of the 4th International Workshop on
%% Mining Software Repositories for Privacy and Security (MSR4P\&S '22)}.
%% ACM, pp.~[insert pages].

%% \bibitem[Wu et al.(2023)]{wu2023autogen}
%% Wu, Q., Bansal, G., Zhang, J., Wu, Y., Li, B., Zhu, E., ... and
%% Wang, C., 2023.
%% \newblock AutoGen: Enabling next-generation LLM applications via
%% multi-agent conversation.
%% \newblock \emph{arXiv preprint arXiv:2308.08155}.

%% \bibitem[Yao et al.(2023)]{yao2023}
%% Yao, S., Zhao, J., Yu, D., Du, N., Shafran, I., Narasimhan, K.
%% and Cao, Y., 2023.
%% \newblock ReAct: Synergizing reasoning and acting in language models.
%% \newblock In: \emph{Proceedings of the 11th International Conference on
%% Learning Representations (ICLR 2023)}.

%% Requires the same preamble as introduction.tex
%% ============================================================

\section{Literature Review}

\subsection{AI-Assisted Programming}

The emergence of large language models trained on code has transformed the ambition of automated programming from the generation of short snippets to the synthesis of complete functional programs. \citet{chen2021} introduced the HumanEval benchmark and the \textit{pass@k} evaluation metric, establishing functional correctness as the field's primary measure of progress. By framing code generation as a problem of producing programs that pass unit tests, HumanEval implicitly defined capability as the ability to satisfy a functional specification; security properties received no representation in the benchmark. The influence of this design choice on subsequent research has been substantial. \citet{li2022alphacode} evaluated AlphaCode on competitive programming contests, achieving [Li et al. 2022 -- insert percentile ranking], a result celebrated as evidence of near-human-level programming ability. \citet{roziere2023} presented Code Llama, demonstrating [Roziere et al. 2023 -- insert benchmark result on HumanEval or equivalent], again evaluated primarily against correctness benchmarks. Across this lineage of models, the optimisation target has been consistent: produce code that executes and passes tests.

The consequences of this framing have been documented empirically. \citet{liu2023} demonstrated that code generated by large language models and evaluated as functionally correct by standard benchmarks nonetheless contains security vulnerabilities at a measurable rate, revealing a systematic gap between the property being measured and the property that matters in deployment. Correctness and security are not the same objective; a function can pass every test in a suite while simultaneously being exploitable if the test suite does not probe security-relevant inputs. The productivity gains that have driven commercial adoption are real: \citet{peng2023} found that developers using AI coding assistance completed tasks [Peng et al. 2023 -- insert productivity improvement, e.g. percentage faster] faster than those working without it. The practical implication is that AI-assisted code is being shipped at an accelerating rate, under the assumption that faster code generation is an unambiguous benefit, without adequate attention to what properties the generated code does and does not possess. The field has, in effect, solved the wrong problem at scale.

\subsection{Security of AI-Generated Code}

The question of whether AI-generated code is secure has attracted a growing body of empirical work, and the findings converge on a troubling pattern that grows more complex as the body of evidence accumulates. \citet{pearce2022} conducted the first large-scale systematic evaluation, generating 1,689 programs across a range of scenarios drawn from the Common Weakness Enumeration taxonomy and finding that approximately 40 per cent contained exploitable security vulnerabilities. Critically, this vulnerability rate was not limited to a subset of task types; it was distributed across diverse CWE categories, establishing that insecure code generation is a property of the model's general behaviour rather than a domain-specific failure mode.

\citet{khoury2023} extended this finding in a particularly revealing direction by demonstrating that ChatGPT produces insecure code even in circumstances where the model itself, when prompted directly, can identify the security principle being violated. This dissociation between declarative security knowledge and generative security practice suggests that the problem is not one of ignorance but of application: the model possesses relevant knowledge but does not activate it during code generation without explicit instruction to do so.

The most consequential finding in this literature, however, comes from \citet{perry2023}, who examined not just the properties of AI-generated code but the effect of AI assistance on developer behaviour. Their controlled study found that participants using AI assistance wrote less secure code than unassisted participants on four out of five tasks, including a particularly striking result on an ECDSA signing task where the proportion of participants producing a secure implementation dropped from 21 per cent in the unassisted condition to 3 per cent in the AI-assisted condition, a difference that was statistically significant ($p = 0.039$). More troubling still, participants in the AI-assisted condition rated their own code as significantly more secure than those in the control condition, despite their outputs being objectively less secure. The confidence paradox documented by Perry et al. constitutes a distinct and underappreciated dimension of the risk: AI assistance does not merely fail to improve security, it actively degrades the developer's epistemic posture toward security.

The breadth of the vulnerability problem across CWE categories has been characterised more formally by \citet{siddiq2022}, whose SecurityEval benchmark provides 130 security-specific samples spanning 75 CWE types, offering a more structured evaluation instrument than task-based studies permit. Taken together, the findings of Pearce et al., Khoury et al., Perry et al., and Siddiq and Santos trace three distinct but related failure modes: AI generates vulnerable code, AI-assisted developers are less likely to notice or correct that vulnerability, and the developer's confidence in the insecure output is, if anything, elevated.

\citet{shukla2025} introduce a fourth failure mode that carries particular implications for iterative AI-assisted workflows. Their study of LLM code refinement found that repeated cycles of AI feedback and revision produced a 37.6 per cent increase in critical vulnerabilities after five iterations. Rather than converging on more secure outputs through successive refinement, the iterative process amplified security weaknesses. This finding challenges the intuitive assumption that more AI involvement in the development process produces better outcomes and suggests that the relationship between the quantity of AI-generated content and code security is not merely neutral but potentially negative.

\subsection{Multi-Agent Systems and Structured Reasoning}

The limitations of single-agent LLM systems on complex tasks have motivated a parallel body of research into multi-agent architectures, in which the division of cognitive labour across specialised agents is proposed as a mechanism for improving output quality. \citet{li2023camel} introduced the CAMEL framework, demonstrating that role-based decomposition, in which agents adopt defined functional identities within a collaborative workflow, improved performance on tasks requiring extended reasoning relative to single-agent baselines. \citet{wu2023autogen} developed AutoGen, a framework for multi-agent conversational systems that enables flexible orchestration of agent roles, showing that conversation-based multi-agent coordination consistently outperformed single-agent approaches on complex compositional tasks.

The theoretical basis for these improvements was partly anticipated by \citet{yao2023}, whose ReAct framework demonstrated that interleaving explicit reasoning traces with action selection produced more reliable task completion than either pure reasoning or pure action. The relevance of ReAct to agentic security pipelines is structural: if explicit reasoning prior to action improves outcome quality, then a pipeline architecture that mandates explicit threat analysis prior to code generation replicates this mechanism at the level of agent orchestration rather than individual token prediction.

The most directly relevant empirical evidence for the deployment of multi-agent human-in-the-loop systems comes from \citet{takerngsaksiri2025}, who documented the HULA system developed and deployed at Atlassian. Across a production workload, HULA achieved plan success in 79 per cent of cases, code success in 87 per cent of cases, and a 59 per cent pull request merge rate, with 568 tasks terminated early by human reviewers exercising their oversight role. These figures represent the most rigorous published evidence for the real-world viability of human-orchestrated agentic pipelines in a software development context. However, the HULA evaluation focused exclusively on task completion and development efficiency. Security outcomes were not measured, vulnerability density was not reported, and the relationship between human oversight and the security properties of the produced code was not investigated. HULA thus demonstrates that the architectural category is viable and scalable; it does not establish whether that architecture produces more secure code than a simpler alternative.

\subsection{Human-AI Collaboration}

The question of how to design effective human oversight of AI systems predates the current generation of code generation tools and has produced a body of design guidance that the AI-assisted programming literature has been slow to engage with. \citet{amershi2019} synthesised 18 guidelines for human-AI interaction from human-computer interaction research, establishing transparency, controllability, and calibrated confidence as foundational requirements for systems that integrate human and AI decision-making. Calibration is particularly relevant in the context of AI-assisted code generation: Perry et al.'s confidence paradox indicates that current tools violate the calibration principle systematically, producing higher confidence in developers when confidence is least warranted.

\citet{shneiderman2020} challenged the prevailing assumption that high automation and high human control represent opposite ends of a spectrum, arguing instead that effective human-centred AI systems achieve both simultaneously through appropriate interface and workflow design. The implication is that reducing AI automation to increase human control is not the correct response to the risks documented by Perry et al. and Shukla et al.; the correct response is to redesign the workflow so that human oversight is structurally effective rather than structurally bypassed.

The evidence on whether human oversight is structurally effective in practice is not encouraging. \citet{perry2023} observed that existing AI coding tools position the developer primarily as a passive acceptor of suggestions, with the cognitive effort required to evaluate a suggestion for security properties working in tension with the productivity incentive that motivates AI adoption in the first place. \citet{kennedy2025} examined human oversight in AI governance contexts more broadly, finding that humans embedded in oversight loops provided correct oversight only approximately half the time, with failure driven predominantly by institutional and social pressure to approve rather than to scrutinise. This finding is significant because it demonstrates that the mere presence of a human decision point in an AI workflow does not guarantee meaningful oversight; the design of that decision point, including the information provided, the framing of the choice, and the presence or absence of structured evaluation criteria, determines whether oversight is genuine or nominal.

\subsection{Research Gap}

The four bodies of literature reviewed above identify a convergent problem that none of them individually resolves. The AI-assisted programming literature has established that current models optimise for correctness and achieve it at scale, but correctness and security are independent properties. The security literature has established that AI-generated code is insecure at a substantial rate, that developer vigilance is reduced by AI assistance, and that iterative refinement compounds rather than corrects these deficiencies. The multi-agent literature has established that structured role-based architectures improve reasoning quality on complex tasks and has demonstrated their viability in production software development, but has not examined security outcomes. The human-AI collaboration literature has established the design principles for effective oversight but has also documented the tendency of oversight mechanisms to become nominal in practice without structural support for genuine scrutiny.

What remains absent from the empirical record is a controlled experiment that brings these threads together: a study that directly compares a standard single-agent code generation baseline against a structured multi-agent pipeline incorporating pre-generation threat modelling, independent review, and formal verification, with security outcome as the primary dependent variable and developer reasoning quality as an explicit secondary outcome. No existing study has reported vulnerability density as a comparative metric across these two architectural conditions under equivalent task and participant constraints. The present study is designed to fill that gap. By holding the underlying model, temperature, and task set constant across conditions, and varying only the architectural structure through which code generation proceeds, the study isolates the effect of multi-agent orchestration on security outcomes in a way that existing work does not permit. The explicit measurement of participant decisions at human-in-the-loop checkpoints further distinguishes this study from both the single-agent vulnerability literature and the multi-agent task-completion literature by treating developer security reasoning as an outcome in its own right rather than an incidental observation.

%% ============================================================
%% BIBLIOGRAPHY ENTRIES TO ADD TO YOUR .bib FILE
%% Add these to the existing entries from introduction.tex
%% ============================================================

%% \bibitem[Amershi et al.(2019)]{amershi2019}
%% Amershi, S., Weld, D., Vorvoreanu, M., Fourney, A., Nushi, B.,
%% Collisson, P., Suh, J., Iqbal, S., Bennett, P., Inkpen, K., Teevan, J.,
%% Kikin-Gil, R. and Horvitz, E., 2019.
%% Software engineering for machine learning: A case study.
%% OR: Guidelines for Human-AI Interaction [verify full title from source].
%% \newblock In: \emph{Proceedings of the 2019 CHI Conference on Human Factors
%% in Computing Systems (CHI '19)}, Glasgow, UK. ACM, pp.~[insert pages].

%% \bibitem[Chen et al.(2021)]{chen2021}
%% Chen, M., Tworek, J., Jun, H., Yuan, Q., Pinto, H.P.O., Kaplan, J., ...
%% and Zaremba, W., 2021.
%% \newblock Evaluating large language models trained on code.
%% \newblock \emph{arXiv preprint arXiv:2107.03374}.

%% \bibitem[Khoury et al.(2023)]{khoury2023}
%% Khoury, R., Avila, A.R., Brunelle, J. and Camara, B.M., 2023.
%% \newblock How secure is code generated by ChatGPT?
%% \newblock In: \emph{[Conference -- insert from paper]}, [Location].
%% IEEE/ACM, pp.~[insert pages].

%% \bibitem[Li et al.(2022)]{li2022alphacode}
%% Li, Y., Choi, D., Chung, J., Kushman, N., Schrittwieser, J., Leblond, R.,
%% ... and Vinyals, O., 2022.
%% \newblock Competition-level code generation with AlphaCode.
%% \newblock \emph{Science}, 378(6624), pp.~1092--1097.

%% \bibitem[Li et al.(2023)]{li2023camel}
%% Li, G., Hammoud, H.A.A.K., Itani, H., Khizbullin, D. and
%% Ghanem, B., 2023.
%% \newblock CAMEL: Communicative agents for ``mind'' exploration of large
%% language model society.
%% \newblock In: \emph{Advances in Neural Information Processing Systems (NeurIPS 2023)},
%% [Location]. NeurIPS, pp.~[insert pages].

%% \bibitem[Liu et al.(2023)]{liu2023}
%% Liu, J. et al., 2023.
%% \newblock [Full title -- verify from paper; relates to security errors in
%% functionally correct AI-generated code].
%% \newblock \emph{[Venue -- insert from paper]}, pp.~[pages].

%% \bibitem[Peng et al.(2023)]{peng2023}
%% Peng, S., Kalliamvakou, E., Croft, P. and Lombard, C., 2023.
%% \newblock The impact of AI on developer productivity: Evidence from
%% GitHub Copilot.
%% \newblock \emph{arXiv preprint arXiv:2302.06590}.

%% \bibitem[Roziere et al.(2023)]{roziere2023}
%% Roziere, B., Gehring, J., Gloeckle, F., Sootla, S., Gat, I., Alon, U.,
%% ... and Synnaeve, G., 2023.
%% \newblock Code Llama: Open foundation models for code.
%% \newblock \emph{arXiv preprint arXiv:2308.12950}.

%% \bibitem[Shneiderman(2020)]{shneiderman2020}
%% Shneiderman, B., 2020.
%% \newblock Human-centered artificial intelligence: Three fresh ideas.
%% \newblock \emph{AIS Transactions on Human-Computer Interaction}, 12(3),
%% pp.~109--124.

%% \bibitem[Siddiq and Santos(2022)]{siddiq2022}
%% Siddiq, M.L. and Santos, J.C.S., 2022.
%% \newblock SecurityEval dataset: Mining vulnerability examples to evaluate
%% machine learning-based code generation techniques.
%% \newblock In: \emph{Proceedings of the 4th International Workshop on
%% Mining Software Repositories for Privacy and Security (MSR4P\&S '22)}.
%% ACM, pp.~[insert pages].

%% \bibitem[Wu et al.(2023)]{wu2023autogen}
%% Wu, Q., Bansal, G., Zhang, J., Wu, Y., Li, B., Zhu, E., ... and
%% Wang, C., 2023.
%% \newblock AutoGen: Enabling next-generation LLM applications via
%% multi-agent conversation.
%% \newblock \emph{arXiv preprint arXiv:2308.08155}.

%% \bibitem[Yao et al.(2023)]{yao2023}
%% Yao, S., Zhao, J., Yu, D., Du, N., Shafran, I., Narasimhan, K.
%% and Cao, Y., 2023.
%% \newblock ReAct: Synergizing reasoning and acting in language models.
%% \newblock In: \emph{Proceedings of the 11th International Conference on
%% Learning Representations (ICLR 2023)}.

%% Requires the same preamble as introduction.tex
%% ============================================================

\section{Literature Review}

\subsection{AI-Assisted Programming}

The emergence of large language models trained on code has transformed the ambition of automated programming from the generation of short snippets to the synthesis of complete functional programs. \citet{chen2021} introduced the HumanEval benchmark and the \textit{pass@k} evaluation metric, establishing functional correctness as the field's primary measure of progress. By framing code generation as a problem of producing programs that pass unit tests, HumanEval implicitly defined capability as the ability to satisfy a functional specification; security properties received no representation in the benchmark. The influence of this design choice on subsequent research has been substantial. \citet{li2022alphacode} evaluated AlphaCode on competitive programming contests, achieving [Li et al. 2022 -- insert percentile ranking], a result celebrated as evidence of near-human-level programming ability. \citet{roziere2023} presented Code Llama, demonstrating [Roziere et al. 2023 -- insert benchmark result on HumanEval or equivalent], again evaluated primarily against correctness benchmarks. Across this lineage of models, the optimisation target has been consistent: produce code that executes and passes tests.

The consequences of this framing have been documented empirically. \citet{liu2023} demonstrated that code generated by large language models and evaluated as functionally correct by standard benchmarks nonetheless contains security vulnerabilities at a measurable rate, revealing a systematic gap between the property being measured and the property that matters in deployment. Correctness and security are not the same objective; a function can pass every test in a suite while simultaneously being exploitable if the test suite does not probe security-relevant inputs. The productivity gains that have driven commercial adoption are real: \citet{peng2023} found that developers using AI coding assistance completed tasks [Peng et al. 2023 -- insert productivity improvement, e.g. percentage faster] faster than those working without it. The practical implication is that AI-assisted code is being shipped at an accelerating rate, under the assumption that faster code generation is an unambiguous benefit, without adequate attention to what properties the generated code does and does not possess. The field has, in effect, solved the wrong problem at scale.

\subsection{Security of AI-Generated Code}

The question of whether AI-generated code is secure has attracted a growing body of empirical work, and the findings converge on a troubling pattern that grows more complex as the body of evidence accumulates. \citet{pearce2022} conducted the first large-scale systematic evaluation, generating 1,689 programs across a range of scenarios drawn from the Common Weakness Enumeration taxonomy and finding that approximately 40 per cent contained exploitable security vulnerabilities. Critically, this vulnerability rate was not limited to a subset of task types; it was distributed across diverse CWE categories, establishing that insecure code generation is a property of the model's general behaviour rather than a domain-specific failure mode.

\citet{khoury2023} extended this finding in a particularly revealing direction by demonstrating that ChatGPT produces insecure code even in circumstances where the model itself, when prompted directly, can identify the security principle being violated. This dissociation between declarative security knowledge and generative security practice suggests that the problem is not one of ignorance but of application: the model possesses relevant knowledge but does not activate it during code generation without explicit instruction to do so.

The most consequential finding in this literature, however, comes from \citet{perry2023}, who examined not just the properties of AI-generated code but the effect of AI assistance on developer behaviour. Their controlled study found that participants using AI assistance wrote less secure code than unassisted participants on four out of five tasks, including a particularly striking result on an ECDSA signing task where the proportion of participants producing a secure implementation dropped from 21 per cent in the unassisted condition to 3 per cent in the AI-assisted condition, a difference that was statistically significant ($p = 0.039$). More troubling still, participants in the AI-assisted condition rated their own code as significantly more secure than those in the control condition, despite their outputs being objectively less secure. The confidence paradox documented by Perry et al. constitutes a distinct and underappreciated dimension of the risk: AI assistance does not merely fail to improve security, it actively degrades the developer's epistemic posture toward security.

The breadth of the vulnerability problem across CWE categories has been characterised more formally by \citet{siddiq2022}, whose SecurityEval benchmark provides 130 security-specific samples spanning 75 CWE types, offering a more structured evaluation instrument than task-based studies permit. Taken together, the findings of Pearce et al., Khoury et al., Perry et al., and Siddiq and Santos trace three distinct but related failure modes: AI generates vulnerable code, AI-assisted developers are less likely to notice or correct that vulnerability, and the developer's confidence in the insecure output is, if anything, elevated.

\citet{shukla2025} introduce a fourth failure mode that carries particular implications for iterative AI-assisted workflows. Their study of LLM code refinement found that repeated cycles of AI feedback and revision produced a 37.6 per cent increase in critical vulnerabilities after five iterations. Rather than converging on more secure outputs through successive refinement, the iterative process amplified security weaknesses. This finding challenges the intuitive assumption that more AI involvement in the development process produces better outcomes and suggests that the relationship between the quantity of AI-generated content and code security is not merely neutral but potentially negative.

\subsection{Multi-Agent Systems and Structured Reasoning}

The limitations of single-agent LLM systems on complex tasks have motivated a parallel body of research into multi-agent architectures, in which the division of cognitive labour across specialised agents is proposed as a mechanism for improving output quality. \citet{li2023camel} introduced the CAMEL framework, demonstrating that role-based decomposition, in which agents adopt defined functional identities within a collaborative workflow, improved performance on tasks requiring extended reasoning relative to single-agent baselines. \citet{wu2023autogen} developed AutoGen, a framework for multi-agent conversational systems that enables flexible orchestration of agent roles, showing that conversation-based multi-agent coordination consistently outperformed single-agent approaches on complex compositional tasks.

The theoretical basis for these improvements was partly anticipated by \citet{yao2023}, whose ReAct framework demonstrated that interleaving explicit reasoning traces with action selection produced more reliable task completion than either pure reasoning or pure action. The relevance of ReAct to agentic security pipelines is structural: if explicit reasoning prior to action improves outcome quality, then a pipeline architecture that mandates explicit threat analysis prior to code generation replicates this mechanism at the level of agent orchestration rather than individual token prediction.

The most directly relevant empirical evidence for the deployment of multi-agent human-in-the-loop systems comes from \citet{takerngsaksiri2025}, who documented the HULA system developed and deployed at Atlassian. Across a production workload, HULA achieved plan success in 79 per cent of cases, code success in 87 per cent of cases, and a 59 per cent pull request merge rate, with 568 tasks terminated early by human reviewers exercising their oversight role. These figures represent the most rigorous published evidence for the real-world viability of human-orchestrated agentic pipelines in a software development context. However, the HULA evaluation focused exclusively on task completion and development efficiency. Security outcomes were not measured, vulnerability density was not reported, and the relationship between human oversight and the security properties of the produced code was not investigated. HULA thus demonstrates that the architectural category is viable and scalable; it does not establish whether that architecture produces more secure code than a simpler alternative.

\subsection{Human-AI Collaboration}

The question of how to design effective human oversight of AI systems predates the current generation of code generation tools and has produced a body of design guidance that the AI-assisted programming literature has been slow to engage with. \citet{amershi2019} synthesised 18 guidelines for human-AI interaction from human-computer interaction research, establishing transparency, controllability, and calibrated confidence as foundational requirements for systems that integrate human and AI decision-making. Calibration is particularly relevant in the context of AI-assisted code generation: Perry et al.'s confidence paradox indicates that current tools violate the calibration principle systematically, producing higher confidence in developers when confidence is least warranted.

\citet{shneiderman2020} challenged the prevailing assumption that high automation and high human control represent opposite ends of a spectrum, arguing instead that effective human-centred AI systems achieve both simultaneously through appropriate interface and workflow design. The implication is that reducing AI automation to increase human control is not the correct response to the risks documented by Perry et al. and Shukla et al.; the correct response is to redesign the workflow so that human oversight is structurally effective rather than structurally bypassed.

The evidence on whether human oversight is structurally effective in practice is not encouraging. \citet{perry2023} observed that existing AI coding tools position the developer primarily as a passive acceptor of suggestions, with the cognitive effort required to evaluate a suggestion for security properties working in tension with the productivity incentive that motivates AI adoption in the first place. \citet{kennedy2025} examined human oversight in AI governance contexts more broadly, finding that humans embedded in oversight loops provided correct oversight only approximately half the time, with failure driven predominantly by institutional and social pressure to approve rather than to scrutinise. This finding is significant because it demonstrates that the mere presence of a human decision point in an AI workflow does not guarantee meaningful oversight; the design of that decision point, including the information provided, the framing of the choice, and the presence or absence of structured evaluation criteria, determines whether oversight is genuine or nominal.

\subsection{Research Gap}

The four bodies of literature reviewed above identify a convergent problem that none of them individually resolves. The AI-assisted programming literature has established that current models optimise for correctness and achieve it at scale, but correctness and security are independent properties. The security literature has established that AI-generated code is insecure at a substantial rate, that developer vigilance is reduced by AI assistance, and that iterative refinement compounds rather than corrects these deficiencies. The multi-agent literature has established that structured role-based architectures improve reasoning quality on complex tasks and has demonstrated their viability in production software development, but has not examined security outcomes. The human-AI collaboration literature has established the design principles for effective oversight but has also documented the tendency of oversight mechanisms to become nominal in practice without structural support for genuine scrutiny.

What remains absent from the empirical record is a controlled experiment that brings these threads together: a study that directly compares a standard single-agent code generation baseline against a structured multi-agent pipeline incorporating pre-generation threat modelling, independent review, and formal verification, with security outcome as the primary dependent variable and developer reasoning quality as an explicit secondary outcome. No existing study has reported vulnerability density as a comparative metric across these two architectural conditions under equivalent task and participant constraints. The present study is designed to fill that gap. By holding the underlying model, temperature, and task set constant across conditions, and varying only the architectural structure through which code generation proceeds, the study isolates the effect of multi-agent orchestration on security outcomes in a way that existing work does not permit. The explicit measurement of participant decisions at human-in-the-loop checkpoints further distinguishes this study from both the single-agent vulnerability literature and the multi-agent task-completion literature by treating developer security reasoning as an outcome in its own right rather than an incidental observation.

%% ============================================================
%% BIBLIOGRAPHY ENTRIES TO ADD TO YOUR .bib FILE
%% Add these to the existing entries from introduction.tex
%% ============================================================

%% \bibitem[Amershi et al.(2019)]{amershi2019}
%% Amershi, S., Weld, D., Vorvoreanu, M., Fourney, A., Nushi, B.,
%% Collisson, P., Suh, J., Iqbal, S., Bennett, P., Inkpen, K., Teevan, J.,
%% Kikin-Gil, R. and Horvitz, E., 2019.
%% Software engineering for machine learning: A case study.
%% OR: Guidelines for Human-AI Interaction [verify full title from source].
%% \newblock In: \emph{Proceedings of the 2019 CHI Conference on Human Factors
%% in Computing Systems (CHI '19)}, Glasgow, UK. ACM, pp.~[insert pages].

%% \bibitem[Chen et al.(2021)]{chen2021}
%% Chen, M., Tworek, J., Jun, H., Yuan, Q., Pinto, H.P.O., Kaplan, J., ...
%% and Zaremba, W., 2021.
%% \newblock Evaluating large language models trained on code.
%% \newblock \emph{arXiv preprint arXiv:2107.03374}.

%% \bibitem[Khoury et al.(2023)]{khoury2023}
%% Khoury, R., Avila, A.R., Brunelle, J. and Camara, B.M., 2023.
%% \newblock How secure is code generated by ChatGPT?
%% \newblock In: \emph{[Conference -- insert from paper]}, [Location].
%% IEEE/ACM, pp.~[insert pages].

%% \bibitem[Li et al.(2022)]{li2022alphacode}
%% Li, Y., Choi, D., Chung, J., Kushman, N., Schrittwieser, J., Leblond, R.,
%% ... and Vinyals, O., 2022.
%% \newblock Competition-level code generation with AlphaCode.
%% \newblock \emph{Science}, 378(6624), pp.~1092--1097.

%% \bibitem[Li et al.(2023)]{li2023camel}
%% Li, G., Hammoud, H.A.A.K., Itani, H., Khizbullin, D. and
%% Ghanem, B., 2023.
%% \newblock CAMEL: Communicative agents for ``mind'' exploration of large
%% language model society.
%% \newblock In: \emph{Advances in Neural Information Processing Systems (NeurIPS 2023)},
%% [Location]. NeurIPS, pp.~[insert pages].

%% \bibitem[Liu et al.(2023)]{liu2023}
%% Liu, J. et al., 2023.
%% \newblock [Full title -- verify from paper; relates to security errors in
%% functionally correct AI-generated code].
%% \newblock \emph{[Venue -- insert from paper]}, pp.~[pages].

%% \bibitem[Peng et al.(2023)]{peng2023}
%% Peng, S., Kalliamvakou, E., Croft, P. and Lombard, C., 2023.
%% \newblock The impact of AI on developer productivity: Evidence from
%% GitHub Copilot.
%% \newblock \emph{arXiv preprint arXiv:2302.06590}.

%% \bibitem[Roziere et al.(2023)]{roziere2023}
%% Roziere, B., Gehring, J., Gloeckle, F., Sootla, S., Gat, I., Alon, U.,
%% ... and Synnaeve, G., 2023.
%% \newblock Code Llama: Open foundation models for code.
%% \newblock \emph{arXiv preprint arXiv:2308.12950}.

%% \bibitem[Shneiderman(2020)]{shneiderman2020}
%% Shneiderman, B., 2020.
%% \newblock Human-centered artificial intelligence: Three fresh ideas.
%% \newblock \emph{AIS Transactions on Human-Computer Interaction}, 12(3),
%% pp.~109--124.

%% \bibitem[Siddiq and Santos(2022)]{siddiq2022}
%% Siddiq, M.L. and Santos, J.C.S., 2022.
%% \newblock SecurityEval dataset: Mining vulnerability examples to evaluate
%% machine learning-based code generation techniques.
%% \newblock In: \emph{Proceedings of the 4th International Workshop on
%% Mining Software Repositories for Privacy and Security (MSR4P\&S '22)}.
%% ACM, pp.~[insert pages].

%% \bibitem[Wu et al.(2023)]{wu2023autogen}
%% Wu, Q., Bansal, G., Zhang, J., Wu, Y., Li, B., Zhu, E., ... and
%% Wang, C., 2023.
%% \newblock AutoGen: Enabling next-generation LLM applications via
%% multi-agent conversation.
%% \newblock \emph{arXiv preprint arXiv:2308.08155}.

%% \bibitem[Yao et al.(2023)]{yao2023}
%% Yao, S., Zhao, J., Yu, D., Du, N., Shafran, I., Narasimhan, K.
%% and Cao, Y., 2023.
%% \newblock ReAct: Synergizing reasoning and acting in language models.
%% \newblock In: \emph{Proceedings of the 11th International Conference on
%% Learning Representations (ICLR 2023)}.
